\begin{claim}{C.69}{Proved}{Relative disk centralizers for Fuchsian actions of either kind}
Let $\Gamma$ be a torsion-free discrete group of
orientation-preserving M\"obius transformations of
$\mathbb D$, with its actual action on
$\overline{\mathbb D}$. No finite-generation
assumption is imposed. Every based loop in
\[
                 C_{\mathcal D_\partial}(\Gamma)
                                                               \tag{C.69.1}
\]
has a based filling in that same group.
The same assertion holds for the actual action
$k^{-1}\Gamma k$, with any specified
$k\in\operatorname{Homeo}(\overline{\mathbb D})$.

More precisely, let $\mathcal L_\Gamma$ be the
limit set, and retain the actual spaces and
restriction map
\[
 \Omega_\Gamma=S^1\setminus\mathcal L_\Gamma,\qquad
 D_\Gamma=\mathbb D\cup\Omega_\Gamma,\qquad
 \operatorname{res}_\Gamma:
 C_{\mathcal D_\partial}(\Gamma)\longrightarrow
 C_{\operatorname{Homeo}(D_\Gamma,\Omega_\Gamma)}(\Gamma).
                                                               \tag{C.69.2}
\]
The notation on the right means pointwise fixing
$\Omega_\Gamma$. The map
$\operatorname{res}_\Gamma$ is an isomorphism
of topological groups with their compact-open
and inverse topologies. This isomorphism
extends the full filling, including its inverse
family, across the original limit set.

The proof also gives the following exact relative
surface comparison. For every second-countable
surface $N$ with boundary, inclusion induces
an injection
\[
 \pi_1\bigl(\operatorname{Homeo}(N,\partial N),\id\bigr)
       \longrightarrow
 \pi_1\bigl(\operatorname{Homeo}(N),\id\bigr).
                                                               \tag{C.69.3}
\]
The proof supplies an explicit correction of an
ambient parameter-disk filling on a proper collar.
It permits infinitely many boundary components
and preserves the original parameter-boundary
loop exactly.

Through the actual spherical quotient and
operator maps of C.68, (C.69.1) gives a relative
equivariant filling on a locally connected
simply connected domain whenever the original
prime-end action is the specified closed-disk
conjugate $k^{-1}\Gamma k$. Its original
restriction kernel is retained. In the plane
application the specified infinity lies outside
the domain and is fixed throughout.
\end{claim}

\begin{proof}
\emph{The original partial compactification and its proper action.}
The group $\Gamma=1$ has
$D_\Gamma=\overline{\mathbb D}$, and (C.69.1)
is the Alexander assertion in C.67.
We first record facts that also cover every
nontrivial $\Gamma$.

For any sequence of distinct elements
$\gamma_n\in\Gamma$, a subsequence has points
$a,b\in\mathcal L_\Gamma$ with
\[
 \gamma_n(z)\longrightarrow a
 \quad\hbox{uniformly on compact subsets of }
                        \overline{\mathbb D}\setminus\{b\}.
                                                               \tag{C.69.4}
\]
Here is the explicit disk calculation.
Write
\[
 \gamma_n(z)=e^{i\theta_n}
                  \frac{z-b_n}{1-\overline b_n z},
                    \qquad b_n=\gamma_n^{-1}(0).
                                                               \tag{C.69.5}
\]
Discreteness prevents a sequence of distinct
$\gamma_n$ from remaining in a compact subset
of the M\"obius group. Taking a subsequence
gives $b_n\to b\in S^1$ and
$e^{i\theta_n}\to e^{i\theta}$.
For $z\ne b$, the numerator at the limit
equals $-b(1-\overline b z)$.
Thus (C.69.4) follows with
$a=-e^{i\theta}b$; the denominator is bounded
away from zero on each stated compact set.
Both $a$ and $b$ are in the limit set,
by applying the definition to
$\gamma_n(0)$ and $\gamma_n^{-1}(0)$.

The limit set is $\Gamma$-invariant: applying
an element to a convergent orbit sequence
gives another orbit sequence, and its
inverse gives equality. Thus the displayed
action preserves $D_\Gamma$.
The space $D_\Gamma$ is a second-countable
surface with boundary $\Omega_\Gamma$.
At a point of $\Omega_\Gamma$ use a small
closed-disk boundary chart disjoint from
the closed set $\mathcal L_\Gamma$.
It is contractible: $(z,t)\mapsto(1-t)z$
is an actual contraction in $D_\Gamma$.

The $\Gamma$-action on this space is free
and proper. A nonidentity element is
hyperbolic or parabolic, because a discrete
elliptic element has finite order and
$\Gamma$ is torsion free. All its fixed
points lie in $\mathcal L_\Gamma$, so
it fixes no point of $D_\Gamma$.
For compact $K\subset D_\Gamma$, a sequence
of distinct translates meeting $K$
would contradict (C.69.4): the compact
$K$ avoids the repelling point $b$,
its images converge uniformly to
$a\in\mathcal L_\Gamma$, and $K$ has
positive distance from that closed
limit set. Consequently the exact quotient
\[
                p:D_\Gamma\longrightarrow
                           N_\Gamma=D_\Gamma/\Gamma
                                                               \tag{C.69.6}
\]
is a regular universal covering of
surfaces with boundary, with precisely
the deck group $\Gamma$.
Evenly covered neighborhoods, Hausdorffness,
and second countability follow exactly
as in C.32, using boundary half disks
when required. Its boundary is
$\Omega_\Gamma/\Gamma$, and its interior
is $\mathbb D/\Gamma$.

\emph{A dense set of controlled boundary accesses.}
Let $E_\Gamma$ be the union of
$\Omega_\Gamma$ and the attracting fixed
points of the hyperbolic elements of
$\Gamma$. Then
\[
                          \overline{E_\Gamma}=S^1.
                                                               \tag{C.69.7}
\]
For completeness this also follows from
(C.69.4). If $\Gamma$ contains a
hyperbolic element $\gamma$, use both
of its endpoints, the attracting
points of $\gamma$ and $\gamma^{-1}$.
For $\zeta\in\mathcal L_\Gamma$ choose
distinct $\gamma_n$ with
$\gamma_n(0)\to\zeta$, and pass to
(C.69.4). At least one of the two
chosen endpoints differs from $b$.
Its images under $\gamma_n$ tend to
$\zeta$ and are the attracting points
of the actual conjugates of $\gamma$
or $\gamma^{-1}$.
Thus the hyperbolic fixed points are
dense in the limit set.

If a nontrivial $\Gamma$ has no
hyperbolic element, all its elements
are parabolic. Two with distinct fixed
points would generate a hyperbolic
element: conjugate one to $z\mapsto z+1$
in the upper half plane, write the
other as a parabolic matrix with
nonzero lower-left entry $c$, and
observe that the trace of its product
with the $n$th power of the first
is its original trace plus $nc$.
Its absolute value exceeds two for
some integer $n$.
Thus all fixed points coincide.
After conjugating that point to
infinity, every element is a real
translation. Discreteness makes
this group cyclic, and its limit
set is a single point.
Its $\Omega_\Gamma$ is dense in
$S^1$. The trivial group has
$\Omega_\Gamma=S^1$.
These cases prove (C.69.7).

\emph{The precise cap-side extension.}
Let
\[
             h\in
 C_{\operatorname{Homeo}(D_\Gamma,\Omega_\Gamma)}(\Gamma).
                                                               \tag{C.69.8}
\]
At a point $p\in\Omega_\Gamma$, it and
its inverse already extend continuously
with value $p$.
At an attracting fixed point $p$ of a
hyperbolic $\gamma$, choose a terminal
ray of the hyperbolic axis ending at $p$.
Partition this ray into the actual
translates $\gamma^n A$ of one compact
fundamental segment $A$.
Its image under $h$ is partitioned
into $\gamma^n h(A)$, by commutation.
The compact set $h(A)$ has iterates
converging uniformly to $p$ under
the hyperbolic M\"obius map.
Thus this image ray ends at the
same $p$. The same argument applies
to $h^{-1}$. These are accesses
with the exact original axis labels.

Fix $\zeta\in S^1$ and $\epsilon>0$.
Choose two points $p,q\in E_\Gamma$
on opposite sides of $\zeta$ in
a sufficiently short circular arc
$I$ containing $\zeta$. Choose an
open crosscut $\alpha\subset\mathbb D$
whose closed arc has endpoints
$p$ and $q$, as follows. Near an
endpoint in $\Omega_\Gamma$, use
a short radial tail. Near a
hyperbolic fixed endpoint use
a terminal axis tail as above.
Truncate those two disjoint tails
and join their interior endpoints
by an embedded arc in a small
convex disk cap around $\zeta$.
The tails may be chosen in
disjoint small endpoint neighborhoods;
the joining arc avoids them except
at its two endpoints.
This gives a crosscut whose closure
and the closed boundary arc
$\overline I$ lie in a convex
neighborhood of diameter less than
$\epsilon$.

Let $V$ be the component of
$\mathbb D\setminus\alpha$ adjacent
to $I$. The Jordan curve
$\overline\alpha\cup\overline I$
bounds this cap. Its filled region
lies in the convex hull of that
curve, so
\[
                           \operatorname{diam} V<\epsilon .
                                                               \tag{C.69.9}
\]
The image $\beta=h^{-1}(\alpha)$
has a continuous embedded closure
with the same two endpoints $p,q$,
by the access calculation.
It is therefore another actual
crosscut. Let $U$ be its complementary
component adjacent to the same
boundary arc $I$.

We verify the side, rather than
inferring it from the endpoint set.
Choose $r\in E_\Gamma\cap I$.
There are points $z_n\in\mathbb D$
with both $z_n\to r$ and
$h^{-1}(z_n)\to r$.
For $r\in\Omega_\Gamma$ this follows
from continuity at $r$; for a
hyperbolic fixed point take
$z_n=\gamma^n z_0$ for its attracting
map and use
$h^{-1}(z_n)=\gamma^n h^{-1}(z_0)$.
Eventually $z_n\in V$ and
$h^{-1}(z_n)\in U$.
Since $h$ maps complementary
components of $\beta$ bijectively
to those of $\alpha$, it follows
that
\[
                               h(U)=V .
                                                               \tag{C.69.10}
\]
The set $U\cup I$ is a neighborhood
of $\zeta$ in the closed disk.
Equations (C.69.9)--(C.69.10)
prove that $h(z)\to\zeta$ as
$z\to\zeta$ from the open disk.
The same proof applies to the
inverse. Thus $h$ extends uniquely
to a boundary-fixed homeomorphism
of $\overline{\mathbb D}$.
Commutation with every original
$\gamma$ survives by density.

This proves algebraic surjectivity
of (C.69.2); injectivity follows
from density of the open disk.
Restriction is continuous. Both
groups are closed Polish groups
in their indicated compact-open
homeomorphism groups, since the
commutator equations and pointwise
boundary constraints are closed.
The Polish open orbit-map argument
of C.26/C.30 applied to this
continuous bijective homomorphism
proves continuity of its inverse.
In particular this gives joint
extension for every compact
parameter family and its inverse
family, uniformly on the closed
disk. It is the full topology
comparison, not merely pointwise
extension of the individual maps.

\emph{The relative parameter-disk correction on an arbitrary surface.}
We prove (C.69.3) independently.
Let $f_t$ be a based loop fixing
$\partial N$ pointwise, and let
$B_z$, $z\in D^2$, be a based
ambient filling with that boundary
loop.
Fix a proper closed collar
\[
                     c:\partial N\times[0,1]\hookrightarrow N .
                                                               \tag{C.69.11}
\]
We recall explicitly that such
a collar can be made proper.
A locally finite surface
triangulation with boundary
subcomplex gives a collar by
the adjacent boundary-edge
rectangles and the finite
interval fans at boundary vertices.
Their normal parameters agree
on the successive common edges
and give a product collar.
Start with any such collar
$c_0$ and choose a compatible
proper metric on $N$.
A positive continuous function
$\delta(b)<1/2$ can be chosen
so that
$d(c_0(b,s),b)<1$ for
$0\leq s\leq\delta(b)$:
local continuity gives the
admissible neighborhoods and
a locally finite partition
of unity gives a positive
continuous minorant.
The collar
$c(b,t)=c_0(b,\delta(b)t)$
is proper. Indeed the $b$
coordinates over a compact
target set lie in its closed
unit neighborhood, which is
compact, and the inverse
image is closed in that
compact set times $[0,1]$.

Each boundary component is
a line or a circle. The
connected based family $B_z$
preserves each component and
its orientation. Put
$\psi_z=B_z|_{\partial N}$.
On a line component choose
a real coordinate and denote
its actual increasing map
by $\Psi_z$.
On a circle component, lift
the joint map
$D^2\times\mathbb R\to S^1$
to a real map $\Psi_z$,
normalized by
$\Psi_{z_*}(u)=u$ at the
specified based parameter
$z_*\in\partial D^2$.
The parameter disk is simply
connected, so this is an
actual continuous lift.
It satisfies
\[
 \Psi_z(u+1)=\Psi_z(u)+1,\qquad
             \Psi_z(u)=u\quad(z\in\partial D^2).
                                                               \tag{C.69.12}
\]
For the second equality,
the original boundary loop
fixes the boundary circle;
the integer discrepancy
is continuous on the connected
parameter boundary and is
zero at $z_*$.
Every $\Psi_z$ is increasing.
The same normalized lifting
of $\psi_z^{-1}$ gives its
inverse, by uniqueness.

Choose a continuous function
$\chi:[0,1]\to[0,1]$ with
$\chi(0)=1$ and $\chi(t)=0$
for $t\geq2/3$.
On each real boundary
coordinate put
\[
             E_{z,t}(u)
                  =(1-\chi(t))u+\chi(t)\Psi_z(u).
                                                               \tag{C.69.13}
\]
This is an increasing proper
homeomorphism. On circle
components it commutes with
integer translation and
therefore descends.
The families and their
inverses are continuous:
strict increase is preserved
by the convex combination,
and compact parameter sets
give common properness bounds
from the compact inverse family
of $\Psi_z$.
Extend to $N$ by
\[
 E_z(c(b,t))=c(E_{z,t}(b),t),\qquad
                  E_z=\id\text{ outside the collar}.
                                                               \tag{C.69.14}
\]
The notation in (C.69.14)
uses the just specified
line or circle coordinates.
The formula is identity
near the outer collar edge,
so it pastes to a homeomorphism.

This is also a continuous
family for infinitely many
boundary components.
Properness of the collar
reduces each compact target
set to a compact subset
of $\partial N\times[0,1]$.
Only finitely many components
meet that subset, since
boundary components are
open in the boundary
one-manifold. All formulas
on those components have
the proved compact-parameter
and inverse continuity.
Thus (C.69.14) is continuous
in the original compact-open
and inverse topology.

We have
$E_z|_{\partial N}=\psi_z$
and $E_z=\id$ on the entire
parameter boundary. Consequently
\[
                            B'_z=E_z^{-1}B_z
                                                               \tag{C.69.15}
\]
is a based filling fixing
$\partial N$ pointwise,
with exactly the prescribed
loop $f_t$ on $\partial D^2$.
This proves (C.69.3) without
assuming a boundary-restriction
fibration or a global choice
of circle lifts over its
whole image.

\emph{Filling on the exact quotient and lifting every label.}
Return to a nontrivial
$\Gamma$ and a based loop
$f_t$ in (C.69.1).
Restriction through (C.69.2)
and descent through (C.69.6)
give a based loop $b_t$
of homeomorphisms of
$N_\Gamma$ fixing its
boundary pointwise, with
the exact equation
\[
                            b_t p=p f_t .
                                                               \tag{C.69.16}
\]
Joint continuity and inverse
continuity descend through
the open covering quotient,
as in C.32.

Choose $x\in\mathbb D$ and
$v=p(x)$ in the interior
of $N_\Gamma$.
By \cite[Theorem~1.1 and
Lemma~2.1]{C_Yagasaki2000},
the identity component of
$\operatorname{Homeo}(N_\Gamma,v)$
is contractible.
Here all exceptions are
excluded explicitly:
the only pointed exceptions
for this assertion are the
plane, the compact disk,
the sphere, and the projective
plane. The actual universal
cover (C.69.6) has nontrivial
infinite fundamental group
$\pi_1(N_\Gamma)=\Gamma$
downstairs, so none occurs.
An interior coordinate disk
gives a continuous local
section of evaluation by
homeomorphisms supported in
that disk. Hence evaluation
is a locally trivial bundle
onto the interior, and its
homotopy exact sequence gives
\[
 \pi_1(\operatorname{Homeo}(N_\Gamma),\id)
             \hookrightarrow\pi_1(\operatorname{Int}N_\Gamma,v).
                                                               \tag{C.69.17}
\]
The evaluation loop $b_t(v)$
is null-homotopic: it is
the projection of the closed
loop $f_t(x)$ in the
contractible open disk.
Thus (C.69.17) supplies an
ambient based filling of
$b_t$, and (C.69.15)
corrects it to a filling
$B'_z$ fixing the whole
boundary of $N_\Gamma$.

Lift the actual map
$(x,z)\mapsto B'_z(p(x))$
on $D_\Gamma\times D^2$
through (C.69.6), with the
identity lift at $z_*$.
The product is simply
connected, so the lift
exists and is unique.
Lift the inverse family
with the same normalization.
Their compositions lift
the projection $p$ and
equal identity at $z_*$,
so uniqueness makes them
mutual inverses everywhere.
For every $\gamma\in\Gamma$,
the two compositions of
the lifted family with
$\gamma$ lift the same
map and agree at the
based parameter. They
therefore agree for all
parameters: every deck
label is retained.

For $x\in\Omega_\Gamma$,
the lifted value remains
in the discrete fiber over
$p(x)$, because $B'_z$
fixes that boundary point.
Connectedness of $D^2$
and the identity value
at $z_*$ force it to
remain exactly $x$.
On the parameter boundary,
uniqueness of lifts over
the connected product
$D_\Gamma\times\partial D^2$
recovers the original
$f_t$ exactly.
This gives the desired
continuous filling in
the right side of (C.69.2).
Its continuous inverse
restriction extends the
entire filling to (C.69.1).

The elementary nontrivial
groups require no omitted
case. For a parabolic
generator, the upper-half-plane
coordinate $z\mapsto z+1$
identifies $D_\Gamma$ with
the closed upper half plane
and $N_\Gamma$ with
$S^1\times[0,\infty)$.
For a hyperbolic generator
$z\mapsto\lambda z$,
$\lambda>1$, logarithmic
polar coordinates
$(s,\theta)\in\mathbb R\times[0,\pi]$
give the deck translation
$s\mapsto s+\log\lambda$
and the compact annular
quotient. In both cases
the pointed group used in
(C.69.17) is contractible
by precisely the quoted
exception lists.
The trivial case was
handled at the start.

\emph{The original conjugacy and the frontier transport.}
Conjugation by the specified
closed-disk homeomorphism
$k$ is a topological group
isomorphism carrying the
centralizer of $k^{-1}\Gamma k$
to (C.69.1).
Transport the proved filling
by its inverse. Every
conjugate operator retains
its original label, and
the boundary is fixed
throughout.

For a spherical domain,
compose with the actual
$T_q$ and prime-end maps
of C.68. The exact
conjugation equation
(C.68.12) proves the
claimed ambient commutation,
and the original restriction
kernel acts identically
on the domain throughout.
The complement is fixed.
When the specified infinity
lies in that complement,
the entire filling descends
to the original plane
homeomorphism group with
the topology comparison
already proved in C.66.
\end{proof}

\emph{Remaining scope.}
The original closed-disk
conjugacy is retained in
this theorem. The construction
does not assert that a
uniformizing conjugacy of
an arbitrary interior proper
action extends to its
original prime-end circle.
Nor does it assert uniform
pasting over infinitely many
different domain conjugacies.
Those parts of the general
C.12 extension problem
remain unresolved.
